\documentclass[preprint,12pt]{elsarticle}

\usepackage{amsmath}
\usepackage{amssymb}
\usepackage{graphicx}
\usepackage{booktabs}
\usepackage[hidelinks]{hyperref}

\journal{Imaging Neuroscience}

\begin{document}

\begin{frontmatter}

\title{What does a released average-subject brain encoder actually predict? A noise
ceiling, a parcel-level replication, and three silent failures}

\author[cuea]{Brian Mwai}
\ead{mwaibryn@gmail.com}
\address[cuea]{Department of Natural Sciences, Faculty of Science, The Catholic University of
Eastern Africa, P.O. Box 62157-00200, Nairobi, Kenya}

\begin{abstract}
Released encoding models are increasingly used as instruments: a checkpoint is downloaded,
run on new stimuli, and its predictions treated as a proxy for cortical response. That use
depends on a property the release does not certify, namely that the specific checkpoint, in
the specific configuration its loader selects, predicts held-out cortex at all. A
self-published audit reports that a widely used trimodal encoder, in the average-subject
configuration its loader forces, is \emph{anti-correlated} with cortex at vertex level,
$r = -0.0145$ against an inter-subject ceiling of $+0.0508$. The claim is about sign rather
than magnitude, comes from a company selling the per-subject alternative, and has not been
independently replicated. We replicate it in parcel space on public data. From the Algonauts
2025 responses the leave-one-subject-out noise ceiling for the target episode is
$r_{\text{ceiling}} = 0.1517$, $95\%$ CI $[0.1484, 0.1551]$, over four subjects, all 1000
Schaefer parcels and 592 timepoints. Running the checkpoint over the first 120~s of that
episode, with prediction and recording placed on one clock, gives a group mean $r = +0.0283$
against a ceiling of $0.0959$ on the same 80 scans, $p = 0.048$ against a circular-shift null,
with three of four subjects positive. The checkpoint is weakly positive and short of decisive,
and it does not show the negative sign the audit reports. We also report three withdrawn values
and the defects behind them, selecting a record by position, inferring orientation from axis
length, and correlating a 2~Hz prediction index by index against a 1.49~s recording, because
each is silent and each is reproducible by anyone attempting this measurement.
\end{abstract}

\begin{keyword}
encoding models \sep fMRI \sep noise ceiling \sep replication \sep model auditing
\end{keyword}

\end{frontmatter}

\section{Introduction}

A released encoding model invites a particular kind of use: treat the checkpoint as an
instrument, run it on stimuli its authors never considered, and read its predictions as a
stand-in for measured cortical response. The appeal is obvious, since the alternative is
scanning people. The risk is equally obvious once stated, which is that a technical report
describing an architecture is not a warranty about the weights released under its name, and
that a checkpoint may be released in a configuration whose predictive validity nobody has
established.

For one widely used trimodal encoder \cite{tribe2025} that risk is not hypothetical. A
commercial laboratory has published an audit of the released checkpoint in its
average-subject configuration \cite{sapient2026}, reporting vertex-level $r = -0.0145$
against a measured inter-subject ceiling of $+0.0508$: anti-correlated with real cortex,
across four subjects and all seven Yeo networks \cite{yeo2011}. Taken at
face value, that would mean every downstream study using this checkpoint as an instrument is
reading a signal with no established relationship to cortex.

The audit should not be accepted uncritically. It is self-published by a company selling the
per-subject alternative, it uses four subjects and one stimulus, and it reports magnitudes
near $0.015$ in either direction, so it is a claim about sign rather than about effect size.
Neither should it be dismissed. It is specific, it is falsifiable, and it concerns a
configuration that a released loader selects by default.

This paper replicates the audit's question on public data, at parcel level, and reports the
three silent defects this work met along the way.

\section{The measurement that stands on its own}

An encoder's predictions are meaningless without knowing how much shared signal exists to
predict. The leave-one-subject-out noise ceiling supplies that: each subject's response is
correlated against the mean of the others, giving a lower bound on what any encoder could
achieve, computed from the recordings alone with no model involved.

\subsection{Data and procedure}

We use the publicly released Algonauts 2025 responses \cite{algonauts2025}, four subjects
viewing the same naturalistic episode, parcellated into 1000 Schaefer parcels
\cite{schaefer2018} in MNI152NLin2009cAsym space.
For each subject we correlate the parcel time series against the mean of the remaining three,
per parcel, over time. Undefined correlations stay undefined: a parcel with no variance over
time has no Pearson correlation, and substituting zero would pull the mean toward zero with
invented values. Confidence intervals are percentile bootstrap over subjects with a fixed
seed, so a reported interval reproduces exactly.

\subsection{Result}

For episode \texttt{s01e01a}, the episode a prediction pass would target,
\begin{equation}
r_{\text{ceiling}} = 0.1517, \qquad 95\%\ \text{CI}\ [0.1484,\ 0.1551],
\end{equation}
over four subjects, 1000 of 1000 parcels defined, and 592 timepoints. Per-subject values are
$0.1553$, $0.1548$, $0.1497$ and $0.1470$.

This is larger than the audit's $+0.0508$ by construction rather than by disagreement.
Averaging vertices within a parcel cancels independent noise and raises correlations, so a
parcel-level ceiling and a vertex-level ceiling are not comparable numbers, and conflating
them would manufacture an agreement or a disagreement that is really a change of spatial
unit.

\section{Three silent failures, and why they are reported}
\label{sec:failures}

Three defects in this work changed a reported number without raising an error. Each is
reported because any group repeating the measurement can reproduce it, and each was identified
the same way: replaying the original code path on the same inputs reproduces the original
number exactly, which establishes the cause rather than merely fitting it.

\subsection{The ceiling: the wrong episode, transposed}

An earlier version of the ceiling gave $0.2247$, CI $[0.2147, 0.2363]$, described as covering
``482 parcels''. That value is withdrawn.

\textbf{The episode was selected by position.} The analysis code carried a placeholder for
the stimulus identifier and fell back to the first key in the response file. That key is a
different episode from the one under prediction, and subjects do not agree on which session an
episode is stored under, so positional selection is inconsistent across files as well.

\textbf{The array was transposed.} The correlation routine expects
$(\text{units}, \text{timepoints})$. The recordings are stored
$(\text{timepoints}, \text{parcels})$, and the orientation guard tested whether the first axis
was longer than the second. For a $(482, 1000)$ array that test cannot fire, so every
correlation ran across the parcel axis. The ``482 parcels'' in the withdrawn figure was the
timepoint count. Replaying the original path reproduces $0.2247$, its interval and the count
of 482 exactly.

\subsection{The prediction: two clocks compared index by index}

The first completed prediction pass reported an encoder correlation of $+0.0344$, CI
$[+0.0198, +0.0433]$, against a ceiling of $0.1356$ over 241 timepoints. That value is also
withdrawn. The checkpoint emits one prediction per $0.5$~s of stimulus: its loader cuts each
segment at the model's 2~Hz feature rate. The recordings are sampled once per $1.49$~s
\cite{algonauts2025}. The analysis truncated both arrays to the shorter length and correlated
them sample by sample, so prediction $k$, at stimulus time $0.5k$~s, was scored against scan
$k$, recorded at $1.49k$~s. By the last sample the two were four minutes apart. Replaying that
path on the saved prediction reproduces $+0.0344$ and all four per-subject values exactly.

A synthetic check shows the size of the error. Two series drawn from one signal, one sampled
as the prediction is and one as the recording is, correlate at $r > 0.999$ once aligned and at
$r = -0.03$ when truncated and compared index by index.

Three practices follow, and we state them as recommendations because all three defects are
generic. Select records by name, never by position. Decide array orientation against a known
constant such as the atlas size, never by comparing axis lengths. And place prediction and
recording on one clock, in seconds, before correlating anything: equal array lengths say
nothing about equal sampling.

\section{The prediction pass}
\label{sec:prediction}

\subsection{Procedure}

\begin{enumerate}
\item Load the released checkpoint \cite{tribe2025} in the average-subject configuration its
loader forces, which is the configuration the audit concerns.
\item Predict the response to the opening of episode \texttt{s01e01a}, the episode whose
ceiling is reported above. The encoder is dataloader-bound, at about 17 minutes of GPU time
per 60~s of video on the hardware available, so a window of the episode is predicted rather
than all of it.
\item Project predictions from the encoder's fsaverage5 vertices into the 1000 Schaefer
parcels \cite{schaefer2018}, rather than upsampling the recordings, which would invent
within-parcel structure the data does not have.
\item Map each recorded scan to the stimulus time it responds to. The checkpoint was trained
against BOLD shifted back by 5~s, so a prediction for stimulus time $t$ estimates the scan
recorded at $t + 5$~s. The prediction is interpolated at $1.49j - 5$~s for each scan $j$ whose
stimulus time lies inside the predicted window; scans outside it are dropped, not
extrapolated.
\item Measure the noise ceiling on exactly the scans used for the encoder, so the two numbers
share one denominator.
\item Correlate per parcel over time and report the sign as measured. Nothing is flipped or
absolute-valued.
\end{enumerate}

The 5~s delay is taken from the checkpoint's training configuration and fixed before any
correlation was computed. Other delays are reported as sensitivity and are not selected from.

\subsection{Inference}

A bootstrap over four subjects has too few distinct resamples to carry an interval, so it is
reported but not relied on. The test is a circular-shift null: the aligned prediction is
rotated against the recordings by every shift of at least ten scans, which keeps each series'
autocorrelation intact and breaks only their alignment, and the group mean correlation is
recomputed at each shift. The one-sided $p$ is the share of shifts reaching the observed value,
counting the observed alignment itself.

\subsection{Result on the first 120 seconds}

The first 120~s of the episode give 80 aligned scans, recording indices 4 to 83.

\begin{table}[htbp]
\centering
\begin{tabular}{@{}lr@{}}
\toprule
Encoder, group mean $r$ & $+0.0283$ \\
Per subject (sub-01, 02, 03, 05) & $-0.0048$, $+0.0316$, $+0.0541$, $+0.0322$ \\
Noise ceiling on the same scans & $0.0959$, CI $[0.0885, 0.1025]$ \\
Circular-shift null, 61 shifts & mean $-0.0048$, sd $0.0191$, max $+0.0351$ \\
One-sided $p$ & $0.048$ \\
\bottomrule
\end{tabular}
\caption{The aligned prediction pass over the first 120~s of \texttt{s01e01a}, parcel level.}
\label{tab:window120}
\end{table}

Three of four subjects are positive and one is indistinguishable from zero. Against the
circular-shift null the group mean sits at the margin of chance, $p = 0.048$. Across response
delays of 3 to 7~s the group mean stays positive, from $+0.0239$ to $+0.0377$, so the sign does
not depend on the delay chosen. The ceiling on these 80 scans is lower than the full-episode
$0.1517$ because it is measured on fewer and earlier scans, which is why the two are never
mixed.

Two minutes of film cannot settle the question. What they do establish is that on this window
the released checkpoint is not anti-correlated with the recordings in parcel space: the one
sign the audit predicts is the one the data here do not show.

\section{Conclusion}

The noise ceiling for the target episode is $0.1517$, CI $[0.1484, 0.1551]$, measured from
public data and independent of any model. The aligned prediction pass over the first 120~s
gives a group mean of $+0.0283$ against a ceiling of $0.0959$ on the same scans, at $p = 0.048$
against a circular-shift null: positive, weak, and short of decisive. It does not reproduce the
negative sign the audit reports, and at parcel level it is not comparable with the audit's
vertex-level $-0.0145$. Three withdrawn values and their causes are reported for the same
reason a null is: each was silent, each is reproducible by anyone attempting this
measurement, and none would have announced itself.

\section*{Code and data availability}

The measurement script, the withdrawn code path, and the artifact recording every value
above are public at \url{https://github.com/brn-mwai/monarch}. The ceiling is produced by
\texttt{scripts/paper3\_noise\_ceiling.py}, which takes the episode as an argument, matches
the record by name and asserts orientation against the atlas size. The prediction pass is
\texttt{scripts/paper3\_prediction.py}, which accepts a saved prediction so the analysis can
be rerun without a GPU; the alignment is \texttt{app/services/temporal\_alignment.py}; the
null and the delay sensitivity are \texttt{scripts/paper3\_alignment\_checks.py}.


\section*{CRediT authorship contribution statement}

\textbf{Brian Mwai:} Conceptualization, Methodology, Software, Validation, Formal analysis,
Investigation, Data curation, Writing -- original draft, Writing -- review and editing,
Visualization. The author is the sole contributor to this work.

\section*{Declaration of competing interest}

The author declares no competing financial interests or personal relationships that could
have appeared to influence the work reported in this paper.

\section*{Funding}

This research received no specific grant from any funding agency in the public, commercial or
not-for-profit sectors. Computation was carried out on publicly available free-tier resources.

\section*{Acknowledgements}

The author thanks Dr.\ Songa Mutambi of the Department of Natural Sciences, The Catholic University of
Eastern Africa, for supervising the project of which this work forms part. The supervision was
academic guidance and approval; the author is solely responsible for the content.

\begin{thebibliography}{9}

\bibitem{tribe2025}
S.~d'Ascoli, J.~Rapin, Y.~Benchetrit, H.~Banville, J.-R.~King,
TRIBE: TRImodal Brain Encoder for whole-brain fMRI response prediction, arXiv:2507.22229
(2025).

\bibitem{sapient2026}
The Sapient Company, The Faceless Brain, self-published technical report, 2026,
\url{https://thesapientcompany.com/research}.

\bibitem{algonauts2025}
A.~T.~Gifford, D.~Bersch, M.~St-Laurent, et al.,
The Algonauts Project 2025 Challenge: How the Human Brain Makes Sense of Multimodal Movies,
arXiv:2501.00504 (2025).

\bibitem{schaefer2018}
A.~Schaefer, R.~Kong, E.~M.~Gordon, T.~O.~Laumann, X.-N.~Zuo, A.~J.~Holmes, S.~B.~Eickhoff,
B.~T.~T.~Yeo,
Local-global parcellation of the human cerebral cortex from intrinsic functional connectivity
MRI,
Cereb. Cortex 28 (2018) 3095--3114.

\bibitem{yeo2011}
B.~T.~T.~Yeo et al.,
The organization of the human cerebral cortex estimated by intrinsic functional connectivity,
J. Neurophysiol. 106 (2011) 1125--1165.

\end{thebibliography}

\end{document}
